% !TEX program = xelatex
\documentclass[11pt]{article}

\renewcommand{\baselinestretch}{1.15}

% ============ 页面边距 ============
\usepackage[letterpaper, margin=0.6in]{geometry}

% ============ 颜色系统 ============
\usepackage{xcolor}
\definecolor{nameColor}{RGB}{26, 42, 58}       % 深空蓝（姓名）
\definecolor{sectionColor}{RGB}{44, 62, 107}   % 深海蓝（章节标题）
\definecolor{accentColor}{RGB}{58, 110, 165}   % 科技蓝（链接/高亮）
\definecolor{dateColor}{RGB}{74, 90, 106}      % 铁灰（日期/次要信息）
\definecolor{bodyColor}{RGB}{45, 45, 45}       % 深灰（正文）
\definecolor{lightGray}{RGB}{208, 215, 224}    % 浅灰（分割线）

% ============ 字体配置 ============
\usepackage{fontspec}
\usepackage{xeCJK}

% 全部使用 PingFang SC（中文 + 英文 + 数字）
\setmainfont{PingFang SC}[
    UprightFont = *-Regular,
    BoldFont = *-Semibold,
    ItalicFont = *-Regular,
]

% 中文同样使用 PingFang SC
\setCJKmainfont{PingFang SC}[
    UprightFont = *-Regular,
    BoldFont = *-Semibold,
]

% 等宽字体（代码、URL 等）也统一用 PingFang SC
\setmonofont{PingFang SC}[
    UprightFont = *-Regular,
    BoldFont = *-Semibold,
]

% ============ 间距与格式 ============
\setlength{\parindent}{0pt}
\setlength{\parskip}{5pt}
\usepackage{enumitem}
\setlist{nosep, left=0pt, itemsep=2pt, label=\textbullet}
\usepackage{url}
\usepackage{hyperref}
\hypersetup{
    colorlinks=true,        % 用颜色代替边框
    linkcolor=accentColor,
    urlcolor=accentColor,    % 链接用科技蓝
    citecolor=dateColor,
}

% ============ 自定义命令 ============
\newcommand{\sectionheader}[1]{%
    \vspace{10pt}%
    \noindent\textcolor{sectionColor}{\textbf{\Large #1}}%
    \vspace{2pt}%
    \hrulefill
    \vspace{6pt}%
}

\newcommand{\entry}[5]{%
    \noindent\textcolor{nameColor}{\textbf{#1}} \hfill \textcolor{dateColor}{\textit{#2}} \\
    \textcolor{dateColor}{\textit{#3}} \hfill \textcolor{dateColor}{\textit{#4}} \\
    \textcolor{bodyColor}{#5}%
    \vspace{6pt}%
}

% ============ 页眉禁用 ============
\usepackage{fancyhdr}
\pagestyle{empty}

% ============ 文档开始 ============
\begin{document}

% ------------- 个人信息 -------------
\begin{center}
    {\huge \textcolor{nameColor}{\textbf{谢搏珩}}} \\[6pt]
    \colorbox{dateColor!15}{%
        \textcolor{dateColor}{
            \href{mailto:2094930710@qq.com}{2094930710@qq.com} \quad | \quad
            +86 19116391912 \quad | \quad
            \href{https://github.com/mildclimate}{GitHub: Mildclimate}
        }%
    } \\[2pt]
    \textcolor{accentColor}{\textbf{求职: AI端侧部署 / 前端开发}}

\end{center}

% ------------- 教育经历 -------------
\sectionheader{教育经历}

\entry{天津大学}{2021.09 -- 2025.06}
    {动画专业 · 学士}{天津}
    {主修：动画运动规律、视听语言、动画分镜设计、三维动画制作}

% ------------- 工作经历 -------------
\sectionheader{工作经历}

\entry{好未来}{2025.07 -- 2025.12}
{前端开发工程师（学而思培优 · 增长方向）}{北京}
    {• 主导低代码平台研发与增长活动落地，构建可视化页面搭建工具，支撑业务老师高效配置营销活动与内部运营场景，覆盖B端配置与C端展示全链路。 \\
     • 负责学而思App内“背单词”H5应用开发，处理复杂跨端兼容与性能优化，集成动画引擎与音频能力，实现本地资源包管理，保障高并发活动场景下的流畅体验。 \\
     • 深度参与需求评审与技术方案设计，打通客户端与前端协作链路，推动增长工具与教学场景的融合，有效提升活动迭代效率与用户体验。}

\entry{魔门塔}{2024.07 -- 2025.01}
    {前端开发实习生（工程稳定组 ES）}{苏州}
    {• 负责内部代码合规扫描平台的前端研发，为研发工程师提供代码合规检测、结果可视化与一键修复指引。 \\
     • 基于 WebSocket 实现扫描任务状态实时推送与结果增量展示，支持 10 万级 QPS 数据刷新，保障大规模代码库扫描场景下的交互流畅度。 \\
     • 与后端算法团队紧密配合，设计扫描规则配置界面与触发机制，将算法扫描结果结构化渲染为可筛选、可追溯的合规报告，有效降低人工审查成本。}

% ------------- 项目经历 -------------
\sectionheader{项目经历}

\entry{视觉重建3D App (React Native + Expo + TF.js)}{}
    {个人项目 \href{https://github.com/Mildclimate/ToonCam}{[  https://github.com/Mildclimate/ToonCam   ]}}{}
    {• 基于 TensorFlow.js 实现手机摄像头视频流的实时语义分割与3D场景重建。 \\
     • 集成 Three.js 渲染引擎，实现真实视频与抽象3D模型的上下分屏实时同步显示。 \\
     • 攻克 iPhone 端推理内存瓶颈，利用 tf.tidy() 与 requestAnimationFrame 循环实现稳定运行。}

% \entry{端侧推理框架性能对比工具}{}
%     {个人项目}{}
%     {• 开发跨平台性能测试工具，对比 TFLite / NCNN / MNN 在移动端的推理延迟与功耗。 \\
%      • 使用 Python 自动化收集数据，生成可视化报告，为技术选型提供依据。}

% ------------- 专业技能 -------------
% \sectionheader{专业技能}
% \begin{itemize}
%     \item \textcolor{nameColor}{\textbf{编程语言：}} \textcolor{bodyColor}{JavaScript / TypeScript，Python，C++（基础阅读与调试）}
%     \item \textcolor{nameColor}{\textbf{移动端 \& 前端：}} \textcolor{bodyColor}{React Native，Expo，React}
%     \item \textcolor{nameColor}{\textbf{AI / 端侧推理：}} \textcolor{bodyColor}{TensorFlow.js，TensorFlow Lite，WebGL 后端，模型量化与压缩}
%     \item \textcolor{nameColor}{\textbf{图形学：}} \textcolor{bodyColor}{Three.js，WebGL，3D场景构建与渲染管线基础}
%     \item \textcolor{nameColor}{\textbf{其他：}} \textcolor{bodyColor}{Git，Xcode，Android Studio，Metro}
% \end{itemize}

% ------------- 其他 -------------
\sectionheader{其他}
\begin{itemize}
    \item \textcolor{bodyColor}{开源贡献：为字节可视化库 \texttt{VChart} 提交代码贡献，有参与开源社区协作经验。\href{https://github.com/VisActor/VChart/pull/3787}{[ PR-3787 ]}}
    \item \textcolor{bodyColor}{AI辅助开发：深度使用 Copilot + DeepSeek V4 Pro 进行日常编码，日均 Token 消耗约 1 亿左右。}
    \item \textcolor{bodyColor}{艺术背景：动画专业出身，具备扎实的美术功底与视觉审美能力（高考美术综合分 260/300，色彩单科 96/100）。}
\end{itemize}

\end{document}